\documentclass[11pt,oneside]{article}

\usepackage[a4paper,margin=27mm,headheight=15pt]{geometry}
\usepackage[T1]{fontenc}
\usepackage[utf8]{inputenc}
\IfFileExists{libertinus.sty}{\usepackage{libertinus}}{\usepackage{lmodern}}
\usepackage{microtype}
\usepackage{amsmath,amssymb,amsthm,mathtools,mathrsfs}
\usepackage{bm}
\usepackage{booktabs,longtable,array,tabularx}
\usepackage{graphicx}
\usepackage{pgfplots}
\usepackage{xcolor}
\usepackage{enumitem}
\usepackage{fancyhdr}
\usepackage{titlesec}
\usepackage{listings}
\usepackage{xurl}
\usepackage{hyperref}
\usepackage[nameinlink,noabbrev]{cleveref}

\definecolor{linkblue}{RGB}{25,75,135}
\definecolor{shadegray}{RGB}{246,247,249}
\definecolor{rulegray}{RGB}{195,201,208}
\pgfplotsset{compat=1.18}
\hypersetup{
  colorlinks=true,
  linkcolor=linkblue,
  citecolor=linkblue,
  urlcolor=linkblue,
  pdftitle={A Formula Atlas for the Thue-Morse Sequence},
  pdfsubject={Digit, valuation, Pascal-row, generating, Fourier, and infinite-product representations},
  pdfkeywords={Thue-Morse, Prouhet, binary digit sum, Lucas theorem, Pascal triangle, Mahler equation, Riesz product, Fabius function}
}

\pagestyle{fancy}
\fancyhf{}
\fancyhead[L]{\small Thue--Morse formula atlas}
\fancyhead[R]{\small\nouppercase{\leftmark}}
\fancyfoot[C]{\thepage}
\renewcommand{\headrulewidth}{0.4pt}

\titleformat{\section}{\Large\bfseries}{\thesection}{0.7em}{}
\titleformat{\subsection}{\large\bfseries}{\thesubsection}{0.7em}{}
\titleformat{\subsubsection}{\normalsize\bfseries}{\thesubsubsection}{0.7em}{}
\setcounter{secnumdepth}{3}
\setcounter{tocdepth}{2}
\setlist[itemize]{leftmargin=2em,itemsep=0.25em,topsep=0.4em}
\setlist[enumerate]{leftmargin=2.3em,itemsep=0.25em,topsep=0.4em}
\setlength{\emergencystretch}{2em}
\allowdisplaybreaks

\newtheorem{theorem}{Theorem}[section]
\newtheorem{proposition}[theorem]{Proposition}
\newtheorem{lemma}[theorem]{Lemma}
\newtheorem{corollary}[theorem]{Corollary}
\newtheorem{conjecture}[theorem]{Conjecture}
\theoremstyle{definition}
\newtheorem{definition}[theorem]{Definition}
\newtheorem{algorithm}[theorem]{Algorithm}
\newtheorem{example}[theorem]{Example}
\theoremstyle{remark}
\newtheorem{remark}[theorem]{Remark}
\newtheorem{warning}[theorem]{Warning}

\newcommand{\R}{\mathbb R}
\newcommand{\C}{\mathbb C}
\newcommand{\Q}{\mathbb Q}
\newcommand{\N}{\mathbb N}
\newcommand{\Z}{\mathbb Z}
\newcommand{\Ftwo}{\mathbb F_2}
\newcommand{\e}{\mathrm e}
\newcommand{\ii}{\mathrm i}
\newcommand{\dd}{\,\mathrm d}
\newcommand{\bigO}{\mathcal O}
\newcommand{\smallo}{o}
\newcommand{\sgn}{\operatorname{sgn}}
\newcommand{\sinc}{\operatorname{sinc}}
\newcommand{\wt}{\operatorname{wt}_2}
\newcommand{\vtwo}{v_2}
\newcommand{\bitsum}{s_2}
\newcommand{\eps}{\varepsilon}
\newcommand{\ind}{\mathbf 1}
\newcommand{\floor}[1]{\left\lfloor #1\right\rfloor}
\newcommand{\ceil}[1]{\left\lceil #1\right\rceil}
\newcommand{\res}[2]{\left[#1\right]_{#2}}
\newcommand{\xor}{\mathbin{\oplus}}
\newcommand{\submask}{\preccurlyeq}
\newcommand{\Up}{\operatorname{up}}
\newcommand{\repo}{\nolinkurl{Analysis/FabiusFunction}}
\newcommand{\lean}[1]{%
  \ifmmode\text{\nolinkurl{#1}}\else\nolinkurl{#1}\fi}

\newenvironment{proofidea}{\par\smallskip\noindent\textit{Proof architecture.}\ }{\hfill$\square$\par\smallskip}
\newenvironment{boxedremark}{%
  \par\smallskip\noindent\begin{center}\begin{minipage}{0.94\textwidth}%
  \color{black}\hrule\vspace{0.6em}\small
}{%
  \vspace{0.6em}\hrule\end{minipage}\end{center}\smallskip
}

\lstdefinestyle{pseudo}{
  basicstyle=\ttfamily\small,
  backgroundcolor=\color{shadegray},
  frame=single,
  rulecolor=\color{rulegray},
  columns=fullflexible,
  keepspaces=true,
  showstringspaces=false,
  xleftmargin=0.7em,
  xrightmargin=0.7em,
  aboveskip=0.8em,
  belowskip=0.8em
}

\title{\bfseries A Formula Atlas for the Thue--Morse Sequence\\[0.35em]
\large Digit, valuation, Pascal-row, generating, Fourier, and infinite-product representations}
\author{Companion report to Section~11.8 of\\
\small Vladimir Reshetnikov's \texttt{ProveIt} Fabius-function development}
\date{25 August 2026}

\begin{document}
\maketitle

\begin{abstract}
Section~11.8 of the main Fabius--Rvachev exposition recovers the zero--one
Thue--Morse bit from one signed row of Pascal's triangle by taking an exact
base-two logarithm.  That identity is only one point in a surprisingly dense
network.  This report develops a formula atlas for the same sequence.  It
passes from binary digits to floor sums, factorial and central-binomial
valuations, carries, Lucas congruences, log-free Pascal-row formulas modulo
three, odious and evil enumerators, lacunary generating products, Euler
transforms, algebraic power series over \(\mathbb F_2\), finite exponential
transforms, Prouhet cancellation, ordinary and Walsh Fourier transforms,
Riesz products, autocorrelation, Dirichlet integrals, and infinite products.

Several identities are especially useful as complements to the existing
binomial--logarithm formula.  Among them are
\(\eps_n=(-1)^{v_2\binom{2n}{n}}\), a formula using only the parity of a
base-two valuation; a logarithm-free Pascal-row formula modulo three; a
nonlinear recurrence and a complete Bell-polynomial expression obtained by
logarithmically differentiating the lacunary Euler product; an integer
algebraic lift whose coefficients reduce modulo two to the Thue--Morse bits;
and exact product-to-sum identities in which finite sine products are Fourier
transforms of dyadic Thue--Morse blocks.  The final section separates results
already formalized in \texttt{ProveIt} from natural Lean modules suggested by
this atlas.  Elementary consequences are proved directly; historically
established deeper results are cited rather than presented as claims of
priority.
\end{abstract}

\tableofcontents
\newpage

\section{Conventions and a map of the formulas}
\label{sec:conventions}

For \(n\in\N\), write
\[
  n=\sum_{j\ge 0} b_j(n)2^j,
  \qquad b_j(n)\in\{0,1\},
  \qquad \bitsum(n)=\sum_{j\ge0}b_j(n).
\]
The zero--one and signed Thue--Morse sequences will be denoted by
\begin{equation}
  \tau(n)=\bitsum(n)\bmod 2\in\{0,1\},
  \qquad
  \eps_n=(-1)^{\bitsum(n)}=(-1)^{\tau(n)}=1-2\tau(n).
  \label{eq:basic-definitions}
\end{equation}
Thus
\[
  (\tau(n))_{n\ge0}=0,1,1,0,1,0,0,1,\ldots,
  \qquad
  (\eps_n)_{n\ge0}=1,-1,-1,1,-1,1,1,-1,\ldots .
\]
We use \(\vtwo(m)\) for the exponent of \(2\) in a positive integer \(m\),
\(a\xor b\) for bitwise exclusive-or, and \(k\submask n\) to mean that every
one-bit of \(k\) is also a one-bit of \(n\).  The notation \(\res{x}{m}\)
means the least nonnegative residue of \(x\) modulo \(m\).

The same object appears below in five guises.
\begin{itemize}
\item As a character of the Boolean group of finite binary strings:
      \(\eps_{a\xor b}=\eps_a\eps_b\).
\item As the parity of a two-adic valuation:
      \(\eps_n=(-1)^{\vtwo\binom{2n}{n}}\).
\item As a statistic of row \(n\) of Pascal's triangle: the number of odd
      entries is \(2^{\bitsum(n)}\).
\item As the coefficient sequence of the lacunary product
      \(\prod_{j\ge0}(1-z^{2^j})\).
\item As both a Walsh character and the coefficient vector whose ordinary
      Fourier intensity generates the Thue--Morse Riesz product.
\end{itemize}
These are not parallel coincidences.  Lucas's theorem turns a Pascal row into
the Boolean submask lattice; the finite generating product is its character
polynomial; Kummer's theorem translates carries into valuations; and the
Fourier formulas are evaluations of that same product on the unit circle.

\begin{boxedremark}
\textbf{Proof-status convention.}
The files \lean{ThueMorseBinomialLog.lean},
\lean{ThueMorseGenerating.lean}, and \lean{ThueMorsePrefix.lean} already
formalize the binomial--logarithm identity, the lacunary generating product,
and the iterated-prefix/Prouhet layer.  Other formulas below are proved on
paper but should not be read as already present in Lean.  A dependency-aware
formalization plan appears in \Cref{sec:formalization}.
\end{boxedremark}

\subsection{A sixteen-term comparison table}

The first values show several encodings side by side.  Here
\(C_n=\binom{2n}{n}\) and \(O_n\) is the number of odd entries in Pascal row
\(n\).

\begin{center}
\small
\begin{tabular}{r@{\quad}lrrrrr}
\toprule
\(n\) & binary & \(\bitsum(n)\) & \(\tau(n)\) & \(\eps_n\) & \(\vtwo(C_n)\) & \(O_n\)\\
\midrule
0  & 0000 & 0 & 0 &  1 & 0 & 1\\
1  & 0001 & 1 & 1 & -1 & 1 & 2\\
2  & 0010 & 1 & 1 & -1 & 1 & 2\\
3  & 0011 & 2 & 0 &  1 & 2 & 4\\
4  & 0100 & 1 & 1 & -1 & 1 & 2\\
5  & 0101 & 2 & 0 &  1 & 2 & 4\\
6  & 0110 & 2 & 0 &  1 & 2 & 4\\
7  & 0111 & 3 & 1 & -1 & 3 & 8\\
8  & 1000 & 1 & 1 & -1 & 1 & 2\\
9  & 1001 & 2 & 0 &  1 & 2 & 4\\
10 & 1010 & 2 & 0 &  1 & 2 & 4\\
11 & 1011 & 3 & 1 & -1 & 3 & 8\\
12 & 1100 & 2 & 0 &  1 & 2 & 4\\
13 & 1101 & 3 & 1 & -1 & 3 & 8\\
14 & 1110 & 3 & 1 & -1 & 3 & 8\\
15 & 1111 & 4 & 0 &  1 & 4 & 16\\
\bottomrule
\end{tabular}
\end{center}

\section{Binary, floor, and bitwise formulas}
\label{sec:binary}

\subsection{Recursion, morphism, and dyadic blocks}

Separating the least significant bit gives the basic recurrence
\begin{equation}
  \tau(2n)=\tau(n),\qquad \tau(2n+1)=1-\tau(n),
  \label{eq:tau-recursion}
\end{equation}
\begin{equation}
  \eps_{2n}=\eps_n,\qquad \eps_{2n+1}=-\eps_n.
  \label{eq:eps-recursion}
\end{equation}
Equivalently, the infinite zero--one word is the fixed point beginning in
\(0\) of the morphism \(0\mapsto 01\), \(1\mapsto 10\).  In signed form, a
block is followed by its negative.  If
\[
  \bm\eps^{(m)}=(\eps_0,\ldots,\eps_{2^m-1}),
\]
then
\begin{equation}
  \bm\eps^{(m+1)}=(\bm\eps^{(m)},-\bm\eps^{(m)}),
  \qquad
  \bm\eps^{(m)}=(1,-1)^{\otimes m}.
  \label{eq:tensor-block}
\end{equation}
The tensor formula will reappear as a one-point Walsh spectrum in
\Cref{sec:walsh}.

A slightly more general block law is often the most convenient executable
formula.

\begin{proposition}[Dyadic block multiplicativity]
For \(m,a,b\in\N\) with \(0\le b<2^m\),
\begin{equation}
  \eps_{2^m a+b}=\eps_a\eps_b,
  \qquad
  \tau(2^m a+b)=\tau(a)\xor\tau(b).
  \label{eq:block-multiplicativity}
\end{equation}
\end{proposition}

\begin{proof}
The binary supports of \(2^m a\) and \(b\) are disjoint, so their digit sums
add without carries.  Taking parity gives the zero--one statement, and
exponentiating \((-1)\) gives the signed statement.
\end{proof}

The same argument works for bitwise exclusive-or without any separation
assumption.

\begin{theorem}[Boolean character formula]
For all \(a,b\in\N\),
\begin{equation}
  \tau(a\xor b)=\tau(a)\xor\tau(b),
  \qquad
  \eps_{a\xor b}=\eps_a\eps_b.
  \label{eq:xor-character}
\end{equation}
Thus \(n\mapsto\eps_n\) is the nontrivial character that takes value \(-1\)
on every coordinate generator of the direct sum
\(\bigoplus_{j\ge0}\mathbb F_2\).
\end{theorem}

\subsection{Floor formulas for digits and digit sums}

Each binary digit is a difference of two quotient floors:
\begin{equation}
  b_j(n)=\floor{\frac{n}{2^j}}-2\floor{\frac{n}{2^{j+1}}}.
  \label{eq:bit-floor}
\end{equation}
Summing telescopically after multiplying by the appropriate weights gives the
familiar digit expansion.  Summing without weights yields a less obvious but
very useful formula.

\begin{proposition}[Legendre floor form of the digit sum]
For every \(n\ge0\),
\begin{equation}
  \bitsum(n)=n-\sum_{j\ge1}\floor{\frac{n}{2^j}},
  \label{eq:digit-sum-floor}
\end{equation}
where the sum is finite.  Consequently
\begin{align}
  \tau(n)
    &=\res{n-\displaystyle\sum_{j\ge1}\floor{n/2^j}}{2},
      \label{eq:tau-floor}\intertext{and}
  \eps_n
    &=(-1)^{\displaystyle\sum_{j\ge0}\floor{n/2^j}}
      =\prod_{j\ge0}(-1)^{\floor{n/2^j}}.
      \label{eq:eps-floor-product}
\end{align}
\end{proposition}

\begin{proof}
Insert \eqref{eq:bit-floor} into \(\sum_jb_j(n)\).  The first floor sum starts
at \(j=0\), the second starts at \(j=1\), and the difference is
\(n-\sum_{j\ge1}\floor{n/2^j}\).  For \eqref{eq:eps-floor-product}, adding the
extra \(j=0\) term changes the exponent from \(n-\bitsum(n)\) to
\(2n-\bitsum(n)\), which has the same parity as \(\bitsum(n)\).
\end{proof}

Combining \eqref{eq:bit-floor} with \eqref{eq:basic-definitions} also gives a
fully explicit finite product containing no digit notation:
\begin{equation}
  \eps_n=
  \prod_{j\ge0}
  \left(
    1-2\floor{\frac{n}{2^j}}+4\floor{\frac{n}{2^{j+1}}}
  \right).
  \label{eq:pure-floor-product}
\end{equation}
Every factor is \(1-2b_j(n)\in\{1,-1\}\), and all sufficiently large factors
are \(1\).

\subsection{Square-wave and sinc-product sign formulas}
\label{sec:sinc-sign}

For nonintegral \(x\),
\begin{equation}
  (-1)^{\floor{x}}=\sgn(\sin \pi x).
  \label{eq:square-wave}
\end{equation}
Applying this to \(x=(n+\tfrac12)/2^j\) in
\eqref{eq:eps-floor-product} produces a continuous-looking sign formula.

\begin{theorem}[Rademacher-sine product]
For every \(n\ge0\),
\begin{equation}
  \eps_n=
  \prod_{j\ge0}
  \sgn\!\left(
    \sin\frac{\pi(n+\tfrac12)}{2^j}
  \right).
  \label{eq:sine-sign-product}
\end{equation}
The product is effectively finite, because its factors are positive once
\(2^j>n+\tfrac12\).
\end{theorem}

The same observation identifies the signs of the infinite sinc product that
occurs in the Fourier theory of the Fabius and Rvachev functions.

\begin{corollary}[Analytic sign encoding]
Define
\begin{equation}
  \Phi(x)=\prod_{j=0}^{\infty}
    \sinc\!\left(\frac{\pi x}{2^j}\right).
  \label{eq:Phi-sinc-product}
\end{equation}
If \(N<x<N+1\), then
\begin{equation}
  \sgn \Phi(x)=\eps_N.
  \label{eq:sinc-interval-sign}
\end{equation}
In particular,
\begin{equation}
  \eps_n=\sgn\Phi\!\left(n+\tfrac12\right).
  \label{eq:sinc-half-sample}
\end{equation}
\end{corollary}

\begin{proof}
The denominators in every sinc factor are positive.  For \(N<x<N+1\), no
multiple of \(2^j\) lies strictly between \(N\) and \(x\), so
\(\floor{x/2^j}=\floor{N/2^j}\).  The sign of the product is therefore
\((-1)^{\sum_{j\ge0}\floor{N/2^j}}=\eps_N\).
\end{proof}

\subsection{The ruler transition and an exact summatory function}

Adding one to \(n-1\) erases exactly \(\vtwo(n)\) trailing one-bits and creates
one new one-bit.  Hence
\begin{equation}
  \bitsum(n)=\bitsum(n-1)+1-\vtwo(n),\qquad n\ge1.
  \label{eq:digit-transition}
\end{equation}

\begin{corollary}[Ruler-function transition]
For \(n\ge1\),
\begin{equation}
  \frac{\eps_n}{\eps_{n-1}}
   =-(-1)^{\vtwo(n)},
  \qquad
  \eps_n=-(-1)^{\vtwo(n)}\eps_{n-1}.
  \label{eq:ruler-transition}
\end{equation}
Thus the places where the sign repeats or flips are controlled by the parity
of the binary ruler function \(\vtwo(n)\).
\end{corollary}

Pairing \(\eps_{2j}+\eps_{2j+1}=0\) gives an unusually sharp summatory law.

\begin{theorem}[Exact prefix discrepancy]
Let
\[
  P(N)=\sum_{0\le n<N}\eps_n.
\]
Then
\begin{equation}
  P(N)=
  \begin{cases}
    0,&N\text{ even},\\
    \eps_{(N-1)/2},&N\text{ odd}.
  \end{cases}
  \label{eq:prefix-discrepancy}
\end{equation}
Equivalently,
\begin{equation}
  P(N)=\ind_{N\text{ odd}}\eps_{\floor{N/2}}.
  \label{eq:prefix-indicator}
\end{equation}
Consequently the number of zero bits and one bits below \(N\) are
\begin{equation}
  \#\{n<N:\tau(n)=0\}=\frac{N+P(N)}2,
  \qquad
  \#\{n<N:\tau(n)=1\}=\frac{N-P(N)}2.
  \label{eq:bit-count-prefix}
\end{equation}
\end{theorem}

\begin{proof}
For \(N=2M\), sum the \(M\) cancelling pairs.  For \(N=2M+1\), the same
pairs leave \(\eps_{2M}=\eps_M\).  Formula \eqref{eq:bit-count-prefix} follows
from \(\eps_n=1-2\tau(n)\).
\end{proof}

\section{Two-adic valuation and carry formulas}
\label{sec:valuations}

The floor sum in \eqref{eq:digit-sum-floor} is also Legendre's formula for
the two-adic valuation of a factorial:
\begin{equation}
  \vtwo(n!)=\sum_{j\ge1}\floor{\frac{n}{2^j}}
           =n-\bitsum(n).
  \label{eq:legendre-two}
\end{equation}
This immediately hides the binary expansion inside ordinary arithmetic.

\begin{theorem}[Factorial-valuation formula]
For every \(n\ge0\),
\begin{equation}
  \tau(n)=\res{n-\vtwo(n!)}{2},
  \qquad
  \eps_n=(-1)^{n-\vtwo(n!)}.
  \label{eq:factorial-valuation-formula}
\end{equation}
Since signs only depend on parity, the exponent may equally be written
\(n+\vtwo(n!)\).
\end{theorem}

A cleaner formula removes even the explicit \(n\) by using one central
binomial coefficient.

\begin{theorem}[Central-binomial valuation formula]
For every \(n\ge0\),
\begin{equation}
  \vtwo\!\binom{2n}{n}=\bitsum(n).
  \label{eq:central-binomial-valuation}
\end{equation}
Consequently
\begin{equation}
  \boxed{
  \tau(n)=\res{\vtwo\binom{2n}{n}}{2},
  \qquad
  \eps_n=(-1)^{\vtwo\binom{2n}{n}}.}
  \label{eq:central-binomial-thue-morse}
\end{equation}
Equivalently,
\begin{equation}
  \tau(n)=\sin^2\!\left(
       \frac{\pi}{2}\vtwo\!\binom{2n}{n}
     \right),
  \qquad
  \eps_n=\cos\!\left(
       \pi\vtwo\!\binom{2n}{n}
     \right).
  \label{eq:central-binomial-trig}
\end{equation}
\end{theorem}

\begin{proof}
By \eqref{eq:legendre-two},
\[
 \vtwo\!\binom{2n}{n}
   =\vtwo((2n)!)-2\vtwo(n!)
   =(2n-\bitsum(2n))-2(n-\bitsum(n)).
\]
A left shift does not change the number of one-bits, so
\(\bitsum(2n)=\bitsum(n)\), and the expression reduces to \(\bitsum(n)\).
\end{proof}

This representation is particularly economical: the full value of the
valuation remembers the binary weight, while its parity is exactly the
Thue--Morse bit.

\subsection{Kummer's theorem and addition with carries}

Kummer's theorem says that \(\vtwo\binom{a+b}{a}\) is the number of carries
when \(a\) and \(b\) are added in base two.  In digit-sum form,
\begin{equation}
  \vtwo\!\binom{a+b}{a}
   =\bitsum(a)+\bitsum(b)-\bitsum(a+b).
  \label{eq:kummer-digit-sum}
\end{equation}
The identity is exact, not merely a congruence: every binary carry decreases
the total digit sum by one at its source position and may propagate, with the
net loss counted by Kummer's valuation.

\begin{theorem}[Carry-corrected multiplicativity]
For all \(a,b\ge0\),
\begin{equation}
  \boxed{
  \eps_{a+b}=\eps_a\eps_b
    (-1)^{\vtwo\binom{a+b}{a}}.}
  \label{eq:carry-corrected-sign}
\end{equation}
In zero--one notation,
\begin{equation}
  \tau(a+b)=
  \tau(a)\xor\tau(b)\xor
  \res{\vtwo\binom{a+b}{a}}{2}.
  \label{eq:carry-corrected-bit}
\end{equation}
Thus ordinary addition differs from bitwise exclusive-or exactly by the
parity of the number of carries.
\end{theorem}

\begin{proof}
Let \(c=\vtwo\binom{a+b}{a}\).  Equation
\eqref{eq:kummer-digit-sum} gives
\(\bitsum(a+b)=\bitsum(a)+\bitsum(b)-c\).  Raise \((-1)\) to both sides.
Subtraction and addition are identical modulo two.
\end{proof}

The ruler transition \eqref{eq:ruler-transition} is the special case
\(a=n-1\), \(b=1\), because
\(\binom{n}{n-1}=n\) and \(\eps_1=-1\).
The carry-free case of \eqref{eq:carry-corrected-sign} gives a useful
criterion:
\begin{equation}
  a\mathbin{\&}b=0
  \quad\Longrightarrow\quad
  \eps_{a+b}=\eps_a\eps_b,
  \label{eq:carry-free}
\end{equation}
which contains the dyadic block formula
\eqref{eq:block-multiplicativity}.

\begin{remark}
Formula \eqref{eq:carry-corrected-sign} can be viewed as a cocycle identity.
The character \(\eps\) is exactly multiplicative for the Boolean-group law
\(\xor\); replacing \(\xor\) by ordinary addition introduces the carry
cocycle \((-1)^{\vtwo\binom{a+b}{a}}\).
\end{remark}

\section{Pascal-triangle formulas beyond the logarithm}
\label{sec:pascal}

\subsection{Lucas's theorem and the submask lattice}

Lucas's theorem modulo two gives
\begin{equation}
  \binom{n}{k}\equiv
  \prod_{j\ge0}\binom{b_j(n)}{b_j(k)}\pmod2.
  \label{eq:lucas-binary}
\end{equation}
Every factor is nonzero precisely when \(b_j(k)\le b_j(n)\).  Therefore
\begin{equation}
  \binom{n}{k}\text{ is odd}
  \quad\Longleftrightarrow\quad
  k\submask n.
  \label{eq:odd-submask}
\end{equation}
The odd positions in row \(n\) are not merely equinumerous with subsets of
the one-bits of \(n\); they \emph{are} that Boolean lattice.

Define the parity-support polynomial
\begin{equation}
  Q_n(z)=\sum_{k=0}^{n}
    \ind_{\binom nk\text{ odd}}z^k.
  \label{eq:row-support-polynomial}
\end{equation}
Then the submask description factors it completely:
\begin{equation}
  Q_n(z)=\prod_{j:b_j(n)=1}(1+z^{2^j})
        =\prod_{j\ge0}(1+z^{2^j})^{b_j(n)}.
  \label{eq:row-support-factorization}
\end{equation}
Setting \(z=1\) gives Glaisher's count.

\begin{theorem}[Glaisher count]
The number \(\nu(n)\) of odd entries in row \(n\) of Pascal's triangle is
\begin{equation}
  \nu(n)=\sum_{k=0}^{n}\ind_{\binom nk\text{ odd}}
        =2^{\bitsum(n)}.
  \label{eq:glaisher-count}
\end{equation}
\end{theorem}

This is the structural input to the formula in Section~11.8 of the main
exposition.  Let
\begin{equation}
  \mathcal X(n)=\sum_{k=0}^{n}(-1)^{\binom nk}.
  \label{eq:X-definition}
\end{equation}
Every odd entry contributes \(-1\), every even entry contributes \(+1\), and
there are \(n+1\) entries.  Hence
\begin{equation}
  \mathcal X(n)=n+1-2\nu(n)
               =n+1-2^{\bitsum(n)+1},
  \label{eq:X-exact}
\end{equation}
so that
\begin{equation}
  n+1-\mathcal X(n)=2^{\bitsum(n)+1}.
  \label{eq:X-power-two}
\end{equation}
Taking an exact logarithm recovers \(\bitsum(n)+1\), and then its parity
recovers \(\tau(n)\), exactly as formalized in
\lean{ThueMorseBinomialLog.lean}.

\subsection{A logarithm-free residue formula modulo three}
\label{sec:mod-three}

The logarithm can be removed because powers of two alternate modulo three:
\(2^r\equiv1,2,1,2,\ldots\).  Let \(\res{x}{3}\in\{0,1,2\}\).

\begin{theorem}[Odd-count formula modulo three]
For every \(n\ge0\),
\begin{equation}
  \boxed{
  \tau(n)=\res{\nu(n)}{3}-1,
  \qquad
  \eps_n=3-2\res{\nu(n)}{3}.}
  \label{eq:odd-count-mod3}
\end{equation}
The residue \(\res{\nu(n)}{3}\) is always \(1\) or \(2\), never \(0\).
\end{theorem}

\begin{proof}
By \eqref{eq:glaisher-count}, \(\nu(n)=2^{\bitsum(n)}\).  If
\(\bitsum(n)\) is even, its residue modulo three is \(1\); if it is odd, the
residue is \(2\).  Subtracting one gives \(\tau(n)\), and
\(1-2\tau(n)\) gives the signed formula.
\end{proof}

A form involving the same signs \((-1)^{\binom nk}\) as the existing
binomial--logarithm identity is obtained by omitting the \(k=0\) term.  Put
\begin{equation}
  S(n)=\sum_{k=1}^{n}(-1)^{\binom nk},
  \qquad S(0)=0.
  \label{eq:S-definition}
\end{equation}
Since the omitted term is \((-1)^{\binom n0}=-1\), one has
\begin{equation}
  S(n)=\mathcal X(n)+1=n+2-2^{\bitsum(n)+1}.
  \label{eq:S-exact}
\end{equation}

\begin{corollary}[Binomial residue and trigonometric formulas]
For every \(n\ge0\),
\begin{align}
  \boxed{\tau(n)=\res{2n+S(n)}{3}},
  \label{eq:reshetnikov-mod3}\intertext{and}
  \boxed{
  \tau(n)=\frac{4}{3}\sin^2\!\left(
     \frac{\pi}{3}\bigl(n-S(n)\bigr)
  \right).}
  \label{eq:reshetnikov-sine}
\end{align}
The corresponding signed form is
\begin{equation}
  \eps_n=1-\frac{8}{3}\sin^2\!\left(
     \frac{\pi}{3}\bigl(n-S(n)\bigr)
  \right).
  \label{eq:reshetnikov-signed-sine}
\end{equation}
\end{corollary}

\begin{proof}
From \eqref{eq:S-exact},
\[
  2n+S(n)=3n+2-2^{\bitsum(n)+1}.
\]
Modulo three this is \(0\) when \(\bitsum(n)\) is even and \(1\) when it is
odd.  Also
\[
  n-S(n)=2^{\bitsum(n)+1}-2,
\]
which is congruent to \(0\) or \(2\) modulo three according as
\(\bitsum(n)\) is even or odd.  The squared sine is therefore \(0\) or
\(3/4\), respectively.
\end{proof}

Formula \eqref{eq:reshetnikov-mod3} was proposed in
\cite{reshetnikovNiceFormula}; the derivation above shows that it is the
modulo-three shadow of the same Glaisher count that drives the exact
base-two logarithm.  In that sense, the logarithmic and log-free formulas are
twins rather than unrelated curiosities.

\subsection{A one-parameter family of signed row statistics}

The signed sum \(\mathcal X(n)\) is the specialization \(u=-1\) of a simple
family.  Define
\begin{equation}
  R_u(n)=\sum_{k=0}^{n}u^{\binom nk\bmod2}.
  \label{eq:Ru-definition}
\end{equation}
Since the exponent is \(0\) at an even entry and \(1\) at an odd entry,
\begin{equation}
  \boxed{R_u(n)=n+1+(u-1)2^{\bitsum(n)}.}
  \label{eq:Ru-exact}
\end{equation}
Thus \(R_{-1}(n)=\mathcal X(n)\), \(R_0(n)=n+1-\nu(n)\) counts the even
entries, and any fixed \(u\ne1\) recovers the odd count by
\begin{equation}
  2^{\bitsum(n)}=\frac{R_u(n)-(n+1)}{u-1}.
  \label{eq:Ru-recovery}
\end{equation}
For real \(u>1\), this gives a continuum of exact logarithmic formulas:
\begin{equation}
  \tau(n)=\frac{1-(-1)^{\log_2((R_u(n)-n-1)/(u-1))}}2.
  \label{eq:Ru-log-family}
\end{equation}
The special choice \(u=-1\) is arithmetically attractive because all
summands are signs, but it is not structurally unique.

\subsection{Pascal probes for individual binary digits}

Lucas's theorem with \(k=2^j\) isolates one coordinate:
\begin{equation}
  b_j(n)=\binom{n}{2^j}\bmod2.
  \label{eq:lucas-bit-probe}
\end{equation}
Hence the whole sequence may be written as a product of selected entries from
a single Pascal row:
\begin{equation}
  \boxed{
  \eps_n=(-1)^{\displaystyle\sum_{j:2^j\le n}\binom{n}{2^j}},
  \qquad
  \tau(n)=\res{\displaystyle\sum_{j:2^j\le n}\binom{n}{2^j}}{2}.}
  \label{eq:pascal-bit-probes}
\end{equation}
No valuation or logarithm appears; the price is that one samples several
specially positioned entries rather than aggregating the full row.

\subsection{A multivariate submask identity}
\label{sec:submask-identity}

If \(k\submask n\), subtraction \(n-k\) has no borrows: every one-bit of
\(n\) is assigned either to \(k\) or to \(n-k\).  Therefore
\begin{theorem}[Submask weight enumerator]
For indeterminates \(x,y\),
\begin{equation}
  \sum_{\substack{0\le k\le n\\\binom nk\text{ odd}}}
     x^{\bitsum(k)}y^{\bitsum(n-k)}
   =(x+y)^{\bitsum(n)}.
  \label{eq:submask-weight-enumerator}
\end{equation}
\end{theorem}

\begin{proof}
For each of the \(\bitsum(n)\) occupied bit positions, choose independently
whether that bit is placed in \(k\), contributing \(x\), or in \(n-k\),
contributing \(y\).  Multiplying the local contributions gives the right-hand
side.
\end{proof}

Several formulas fall out by specialization:
\begin{align}
  \sum_{\binom nk\text{ odd}}z^{\bitsum(k)}
    &=(1+z)^{\bitsum(n)},
    \label{eq:submask-one-variable}\\
  \sum_{\binom nk\text{ odd}}\eps_k
    &=(1-1)^{\bitsum(n)}=\ind_{n=0},
    \label{eq:submask-sign-cancel}\\
  \boxed{
  \eps_n
    &=\sum_{\substack{0\le k\le n\\\binom nk\text{ odd}}}
      (-2)^{\bitsum(k)}.}
    \label{eq:submask-minus-two}
\end{align}
The last identity is simply \((1-2)^{\bitsum(n)}\), but written entirely as
a sum over the odd positions of Pascal row \(n\).  Replacing
\(\bitsum(k)\) by \(\vtwo\binom{2k}{k}\) using
\eqref{eq:central-binomial-valuation} yields the nested binomial formula
\begin{equation}
  \eps_n=
  \sum_{k=0}^{n}
  \ind_{\binom nk\text{ odd}}
  (-2)^{\vtwo\binom{2k}{k}}.
  \label{eq:nested-binomial-formula}
\end{equation}
It is not computationally efficient, but it displays a curious closure of the
sequence under two different binomial encodings.

\section{Odious and evil enumerators}
\label{sec:odious}

A nonnegative integer is called \emph{evil} when its binary digit sum is even
and \emph{odious} when it is odd.  Let \(\mathcal E_n\) and
\(\mathcal O_n\) be the \(n\)-th evil and odious numbers, both indexed from
\(n=0\):
\[
  \mathcal E_0,\mathcal E_1,\ldots=0,3,5,6,9,10,12,15,\ldots,
\]
\[
  \mathcal O_0,\mathcal O_1,\ldots=1,2,4,7,8,11,13,14,\ldots .
\]
The two increasing enumerators admit remarkably short formulas
\cite{alloucheCloitreShevelev}.

\begin{theorem}[Evil and odious enumerators]
For every \(n\ge0\),
\begin{equation}
  \boxed{
  \mathcal E_n=2n+\tau(n),
  \qquad
  \mathcal O_n=2n+1-\tau(n).}
  \label{eq:evil-odious-enumerators}
\end{equation}
Consequently
\begin{equation}
  \mathcal O_n-\mathcal E_n=\eps_n,
  \qquad
  \mathcal O_n+\mathcal E_n=4n+1.
  \label{eq:evil-odious-difference}
\end{equation}
Thus the signed Thue--Morse sequence is exactly the termwise displacement
between the two complementary enumerations.
\end{theorem}

\begin{proof}
The two candidates are the consecutive integers \(2n\) and \(2n+1\).  Their
binary digit-sum parities are \(\tau(n)\) and \(1-\tau(n)\), respectively.
If \(\tau(n)=0\), the even candidate is evil and the odd candidate is odious;
if \(\tau(n)=1\), the roles are reversed.  The formulas select the correct
member in each case.  Monotonicity follows because each pair belongs to the
interval \(\{2n,2n+1\}\).
\end{proof}

The formulas are stable under composition because the output parity is known
in advance.

\begin{corollary}[Composition laws]
For every \(n\ge0\),
\begin{align}
  \mathcal E_{\mathcal E_n}&=2\mathcal E_n,
  &\mathcal O_{\mathcal E_n}&=2\mathcal E_n+1,
  \label{eq:EE-OE}\\
  \mathcal E_{\mathcal O_n}&=2\mathcal O_n+1,
  &\mathcal O_{\mathcal O_n}&=2\mathcal O_n.
  \label{eq:EO-OO}
\end{align}
\end{corollary}

\begin{proof}
By definition, \(\tau(\mathcal E_n)=0\) and
\(\tau(\mathcal O_n)=1\).  Substitute those two values into
\eqref{eq:evil-odious-enumerators} with the index replaced by
\(\mathcal E_n\) or \(\mathcal O_n\).
\end{proof}

The prefix discrepancy formula \eqref{eq:prefix-discrepancy} simultaneously
gives exact counting functions:
\begin{equation}
  \#\{\mathcal E_n<N\}=\frac{N+P(N)}2,
  \qquad
  \#\{\mathcal O_n<N\}=\frac{N-P(N)}2.
  \label{eq:evil-odious-counting}
\end{equation}
Their discrepancy never exceeds one.  This is the strongest possible form of
asymptotic equidistribution: every even-length initial interval contains
exactly as many evil as odious numbers.

\section{Generating functions, Euler transforms, and recurrences}
\label{sec:generating}

\subsection{The finite mother product}

For a dyadic block define
\begin{equation}
  P_m(z)=\sum_{n=0}^{2^m-1}\eps_nz^n.
  \label{eq:Pm-definition}
\end{equation}
The block recursion \eqref{eq:tensor-block} gives
\begin{equation}
  P_{m+1}(z)=P_m(z)-z^{2^m}P_m(z)
            =(1-z^{2^m})P_m(z).
  \label{eq:Pm-recursion}
\end{equation}
Starting from \(P_0(z)=1\) yields the fundamental factorization
\begin{equation}
  \boxed{
  P_m(z)=\prod_{j=0}^{m-1}(1-z^{2^j}).}
  \label{eq:finite-generating-product}
\end{equation}
This is formalized coefficientwise in \lean{ThueMorseGenerating.lean}.  It is
the common source of the Prouhet, Fourier, trigonometric, and Riesz-product
formulas developed later.

Coefficientwise stabilization as \(m\to\infty\) gives the formal power series
\begin{equation}
  E(z)=\sum_{n\ge0}\eps_nz^n
      =\prod_{j\ge0}(1-z^{2^j}).
  \label{eq:infinite-signed-generating}
\end{equation}
Analytically, the product converges for \(|z|<1\).  Splitting off its first
factor produces the Mahler equation
\begin{equation}
  \boxed{E(z)=(1-z)E(z^2).}
  \label{eq:E-Mahler}
\end{equation}
For the zero--one generating series
\begin{equation}
  T(z)=\sum_{n\ge0}\tau(n)z^n,
  \label{eq:T-definition}
\end{equation}
we have
\begin{equation}
  T(z)=\frac{1}{2(1-z)}-\frac12E(z),
  \label{eq:T-from-E}
\end{equation}
so the bit recursion gives
\begin{equation}
  \boxed{
  T(z)=(1-z)T(z^2)+\frac{z}{1-z^2}.}
  \label{eq:T-Mahler}
\end{equation}

\subsection{The logarithmic Euler transform}
\label{sec:euler-transform}

Taking the formal logarithm of \eqref{eq:infinite-signed-generating} gives
\begin{align}
  \log E(z)
   &=\sum_{j\ge0}\log(1-z^{2^j})
    =-\sum_{j\ge0}\sum_{r\ge1}\frac{z^{r2^j}}r.
  \label{eq:logE-double}
\end{align}
For a fixed exponent \(k\), the possible \(j\)'s are
\(0\le j\le\vtwo(k)\), with \(r=k/2^j\).  Therefore
\begin{equation}
  \boxed{
  \log E(z)=-\sum_{k\ge1}\frac{a_k}{k}z^k,
  \qquad
  a_k=2^{\vtwo(k)+1}-1.}
  \label{eq:logE-ruler}
\end{equation}
Thus the logarithm of the Thue--Morse coefficient series is controlled by a
simple transform of the binary ruler function.

Logarithmic differentiation yields a nonlinear convolution recurrence that
contains no binary digits.

\begin{theorem}[Ruler-convolution recurrence]
Let \(a_k=2^{\vtwo(k)+1}-1\).  Then \(\eps_0=1\) and, for \(n\ge1\),
\begin{equation}
  \boxed{
  n\eps_n=-\sum_{k=1}^{n}a_k\eps_{n-k}.}
  \label{eq:ruler-convolution}
\end{equation}
\end{theorem}

\begin{proof}
Differentiate \eqref{eq:logE-ruler}:
\[
  \frac{E'(z)}{E(z)}=-\sum_{k\ge1}a_kz^{k-1}.
\]
Multiply by \(E(z)=\sum_{r\ge0}\eps_rz^r\), then compare the coefficient of
\(z^{n-1}\).
\end{proof}

The recurrence is arithmetically surprising: the weighted convolution on the
right is always exactly \(\pm n\).  It also provides a direct route to
complete coefficient formulas by exponentiating \eqref{eq:logE-ruler}.

\begin{theorem}[Partition and Bell-polynomial formulas]
For \(n\ge0\),
\begin{equation}
  \boxed{
  \eps_n=
  \sum_{\substack{m_1,\ldots,m_n\ge0\\
                  m_1+2m_2+\cdots+nm_n=n}}
  \prod_{j=1}^{n}
    \frac{1}{m_j!}
    \left(-\frac{a_j}{j}\right)^{m_j}.}
  \label{eq:partition-formula}
\end{equation}
Equivalently, in terms of the complete exponential Bell polynomial \(B_n\),
\begin{equation}
  \boxed{
  \eps_n=\frac1{n!}
  B_n\bigl(-0!a_1,-1!a_2,\ldots,-(n-1)!a_n\bigr).}
  \label{eq:bell-formula}
\end{equation}
\end{theorem}

\begin{proof}
Equation \eqref{eq:logE-ruler} writes \(E(z)\) as
\[
  \prod_{j\ge1}\exp\!\left(-\frac{a_j}{j}z^j\right).
\]
Expand each exponential and collect terms of total degree \(n\), obtaining
\eqref{eq:partition-formula}.  The Bell-polynomial form is the standard
identity
\[
  \exp\!\left(\sum_{j\ge1}x_j\frac{z^j}{j!}\right)
   =\sum_{n\ge0}B_n(x_1,\ldots,x_n)\frac{z^n}{n!}
\]
with \(x_j=-(j-1)!a_j\).
\end{proof}

The right sides of \eqref{eq:partition-formula} and \eqref{eq:bell-formula}
look rational and highly nonlocal, yet collapse to \(\pm1\).  They are not
recommended as evaluators, but they expose the Thue--Morse sequence as the
Euler transform of the ruler-derived sequence \((a_k)\).

\subsection{An algebraic equation over \(\mathbb F_2\)}
\label{sec:F2-algebraic}

Reduce \(T(z)\) coefficientwise modulo two.  In characteristic two,
\(T(z^2)=T(z)^2\), \(1-z=1+z\), and
\(1+z^2=(1+z)^2\).  Equation \eqref{eq:T-Mahler} becomes
\begin{equation}
  (1+z)^3T(z)^2+(1+z)^2T(z)+z=0
  \qquad\text{in }\Ftwo[[z]].
  \label{eq:T-algebraic-F2}
\end{equation}
Equivalently,
\begin{equation}
  (1+z)T(z)^2+T(z)=\frac{z}{1+z^2}.
  \label{eq:T-algebraic-F2-short}
\end{equation}
This is a concrete instance of Christol's theorem: a sequence over a finite
field is automatic exactly when its generating series is algebraic over the
rational-function field.  Here the equation is obtained directly, with no
appeal to the general theorem.

\subsection{An integer algebraic lift}
\label{sec:integer-lift}

There is a rather different recurrence whose integer values have
Thue--Morse parity.  Define \(a_0=0\), \(a_1=1\), and for \(n\ge2\),
\begin{equation}
  n a_n=(5-2n)a_{n-1}+3(n-1)a_{n-2}+1.
  \label{eq:integer-lift-recurrence}
\end{equation}
The first terms are
\[
  0,1,1,2,1,4,-2,13,-23,68,-164,439,-1146,3067,\ldots .
\]
This recurrence was proposed in \cite{reshetnikovRecurrence}; the formal-power
series proof below follows the answer recorded there.

\begin{theorem}[Algebraic integer lift]
Let \(A(z)=\sum_{n\ge0}a_nz^n\).  Then
\begin{equation}
  A(z)=\frac{1}{2(1-z)}
  \left(\sqrt{\frac{1+3z}{1-z}}-1\right)
  \label{eq:A-closed-form}
\end{equation}
and
\begin{equation}
  (1-z)^3A(z)^2+(1-z)^2A(z)=z.
  \label{eq:A-algebraic}
\end{equation}
All coefficients \(a_n\) are integers, and
\begin{equation}
  \boxed{a_n\equiv\tau(n)\pmod2.}
  \label{eq:A-parity}
\end{equation}
Consequently
\begin{equation}
  \eps_n=(-1)^{a_n}.
  \label{eq:A-sign-formula}
\end{equation}
\end{theorem}

\begin{proof}
Multiplying \eqref{eq:integer-lift-recurrence} by \(z^{n-1}\), summing, and
including the initial terms gives the first-order differential equation
\[
  A'(z)=3A(z)-2zA'(z)+3zA(z)+3z^2A'(z)+\frac1{1-z}.
\]
The solution with \(A(0)=0\) is \eqref{eq:A-closed-form}, and squaring gives
\eqref{eq:A-algebraic}.  Conversely, substitution of a formal series with
zero constant term into \eqref{eq:A-algebraic} determines every successive
coefficient using only integer arithmetic; this gives integrality and
uniqueness.

Reduce \eqref{eq:A-algebraic} modulo two.  It becomes exactly
\eqref{eq:T-algebraic-F2}.  That equation also has a unique solution in
\(z\Ftwo[[z]]\): at each degree, the linear term contributes the new
coefficient with coefficient one, while all other contributions involve
previous coefficients (the Frobenius square only uses even indices).  Hence
\(A(z)\equiv T(z)\pmod2\), proving \eqref{eq:A-parity}.
\end{proof}

\begin{boxedremark}
The lift \((a_n)\) is not a disguised digit sum: it grows in magnitude and
changes sign irregularly.  Its parity is Thue--Morse because an algebraic
equation over \(\mathbb Z[[z]]\) reduces to the characteristic-two equation
of the automatic sequence.  This is a general discovery pattern: search for
an integral algebraic series whose defining polynomial reduces to the known
Christol polynomial.
\end{boxedremark}

\section{Finite exponential transforms and Prouhet cancellation}
\label{sec:exponential}

\subsection{A multivariate mother identity}

The finite generating product is a specialization of a more elementary
Boolean-cube identity.  Let \(n<2^m\) have bits
\(b_0(n),\ldots,b_{m-1}(n)\).

\begin{theorem}[Boolean-cube product]
For arbitrary commuting indeterminates \(x_0,\ldots,x_{m-1}\),
\begin{equation}
  \boxed{
  \sum_{n=0}^{2^m-1}\eps_n
    \prod_{j=0}^{m-1}x_j^{b_j(n)}
  =\prod_{j=0}^{m-1}(1-x_j).}
  \label{eq:boolean-mother}
\end{equation}
\end{theorem}

\begin{proof}
Identify \(n\) with its bit vector \(b\in\{0,1\}^m\).  Then
\(\eps_n=\prod_j(-1)^{b_j}\), and the sum factors into independent
one-coordinate sums:
\[
  \sum_{b\in\{0,1\}^m}\prod_j(-x_j)^{b_j}
   =\prod_j(1-x_j).
\]
\end{proof}

Setting \(x_j=z^{2^j}\) recovers
\eqref{eq:finite-generating-product}.  Setting
\(x_j=\e^{2^jt}\) gives the finite exponential transform
\begin{equation}
  F_m(t)=\sum_{n=0}^{2^m-1}\eps_n\e^{nt}
        =\prod_{j=0}^{m-1}(1-\e^{2^jt}).
  \label{eq:finite-exponential-product}
\end{equation}
Using
\(1-\e^u=-2\e^{u/2}\sinh(u/2)\) and
\(\sum_{j=0}^{m-1}2^j=2^m-1\), one obtains
\begin{equation}
  F_m(t)=(-2)^m
    \e^{(2^m-1)t/2}
    \prod_{j=0}^{m-1}\sinh(2^{j-1}t).
  \label{eq:exponential-sinh}
\end{equation}
After centering the block at \(c_m=(2^m-1)/2\), the exponential factor
disappears:
\begin{equation}
  \boxed{
  \sum_{n=0}^{2^m-1}\eps_n\e^{(n-c_m)t}
   =(-2)^m\prod_{j=0}^{m-1}\sinh(2^{j-1}t).}
  \label{eq:centered-exponential}
\end{equation}
This is the exponential counterpart of the sine-product identity in
\Cref{sec:trig-products}.

\subsection{All finite power sums}

Define the signed power moments of the first dyadic block by
\begin{equation}
  M_{m,r}=\sum_{n=0}^{2^m-1}\eps_n n^r.
  \label{eq:power-moment-definition}
\end{equation}
Since \(F_m(t)=\sum_{r\ge0}M_{m,r}t^r/r!\), the zero of order \(m\) at
\(t=0\) in \eqref{eq:finite-exponential-product} gives Prouhet cancellation.

\begin{theorem}[Prouhet moments and the sharp boundary]
For \(m\ge0\),
\begin{equation}
  M_{m,r}=0\qquad(0\le r<m),
  \label{eq:prouhet-zero}
\end{equation}
while
\begin{equation}
  \boxed{
  M_{m,m}=(-1)^m m!\,2^{\binom m2}.}
  \label{eq:prouhet-boundary}
\end{equation}
\end{theorem}

\begin{proof}
Near zero,
\[
  \prod_{j=0}^{m-1}(1-\e^{2^jt})
  =\prod_{j=0}^{m-1}\bigl(-2^jt+\bigO(t^2)\bigr)
  =(-1)^m2^{0+1+\cdots+(m-1)}t^m+\bigO(t^{m+1}).
\]
Compare exponential-series coefficients.
\end{proof}

The whole range \(r\ge m\) has an explicit composition formula, not merely a
recurrence.

\begin{theorem}[Complete composition formula]
For \(r\ge m\ge1\),
\begin{equation}
  \boxed{
  M_{m,r}=(-1)^m r!\,2^{\binom m2}
  \sum_{\substack{q_0,\ldots,q_{m-1}\ge0\\
                  q_0+\cdots+q_{m-1}=r-m}}
  \prod_{j=0}^{m-1}
    \frac{2^{jq_j}}{(q_j+1)!}.}
  \label{eq:complete-power-moment}
\end{equation}
\end{theorem}

\begin{proof}
Factor the lowest-order term from every factor in
\eqref{eq:finite-exponential-product}:
\begin{align*}
 F_m(t)
 &=(-1)^m2^{\binom m2}t^m
   \prod_{j=0}^{m-1}
   \frac{\e^{2^jt}-1}{2^jt}\\
 &=(-1)^m2^{\binom m2}t^m
   \prod_{j=0}^{m-1}
   \left(\sum_{q\ge0}\frac{2^{jq}t^q}{(q+1)!}\right).
\end{align*}
The coefficient of \(t^r\) is the displayed weak-composition sum, and
multiplication by \(r!\) converts it from an ordinary coefficient to
\(M_{m,r}\).
\end{proof}

Centering adds a parity selection rule.  The right side of
\eqref{eq:centered-exponential} is an even function when \(m\) is even and an
odd function when \(m\) is odd.

\begin{corollary}[Centered parity cancellation]
For \(c_m=(2^m-1)/2\),
\begin{equation}
  \sum_{n=0}^{2^m-1}\eps_n(n-c_m)^r=0
  \quad\text{if }r<m\text{ or }r\not\equiv m\pmod2.
  \label{eq:centered-moment-parity}
\end{equation}
\end{corollary}

The same cancellation applies to every polynomial of degree below \(m\), and
the leading coefficient controls the boundary case.  If \(p\in\Q[x]\) has
degree at most \(m\), then
\begin{equation}
  \sum_{n=0}^{2^m-1}\eps_np(n)
   =[x^m]p(x)\;(-1)^m m!2^{\binom m2}.
  \label{eq:polynomial-coefficient-extraction}
\end{equation}
This coefficient-extraction theorem is formalized in
\lean{ThueMorsePrefix.lean}.

For the binomial basis, whose leading coefficient is \(1/r!\), one gets
\begin{equation}
  \sum_{n=0}^{2^m-1}\eps_n\binom nr=0\quad(r<m),
  \qquad
  \sum_{n=0}^{2^m-1}\eps_n\binom nm=(-1)^m2^{\binom m2}.
  \label{eq:binomial-moments}
\end{equation}

\subsection{Prouhet's equal-sums partition}

Let
\[
  A_m=\{0\le n<2^m:\tau(n)=0\},
  \qquad
  B_m=\{0\le n<2^m:\tau(n)=1\}.
\]
Since \(\eps_n=1\) on \(A_m\) and \(-1\) on \(B_m\),
\eqref{eq:prouhet-zero} is equivalent to
\begin{equation}
  \sum_{n\in A_m}n^r=\sum_{n\in B_m}n^r,
  \qquad 0\le r<m.
  \label{eq:prouhet-partition}
\end{equation}
Thus the first \(2^m\) nonnegative integers split into two classes with equal
sums of like powers through degree \(m-1\), Prouhet's 1851 construction.  At
degree \(m\), their difference is the nonzero quantity in
\eqref{eq:prouhet-boundary}; the order of cancellation is exact.

\section{Ordinary Fourier and trigonometric formulas}
\label{sec:fourier}

\subsection{The finite Fourier product}

Place \(z=\e^{\ii\theta}\) in
\eqref{eq:finite-generating-product}.  Using
\(1-\e^{\ii u}=-2\ii\e^{\ii u/2}\sin(u/2)\) gives
\begin{equation}
  \boxed{
  \sum_{n=0}^{2^m-1}\eps_n\e^{\ii n\theta}
  =(-2\ii)^m
   \e^{\ii(2^m-1)\theta/2}
   \prod_{j=0}^{m-1}\sin(2^{j-1}\theta).}
  \label{eq:finite-Fourier-product}
\end{equation}
Every product-to-sum formula in this section is a real or imaginary part of
this identity after a change of scale.

\subsection{Sine and cosine products as Thue--Morse sums}
\label{sec:trig-products}

Set \(\theta=2z/2^m\), multiply
\eqref{eq:finite-Fourier-product} by
\(\e^{\ii z/2^m}\ii^m\), and simplify the phase.  One obtains the compact
complex identity
\begin{equation}
  \sum_{n=0}^{2^m-1}\eps_n
   \exp\!\left(\ii\left(
      \frac{\pi m}{2}+\frac{(2n+1)z}{2^m}
   \right)\right)
  =2^m\e^{\ii z}\prod_{r=1}^{m}\sin\frac{z}{2^r}.
  \label{eq:complex-product-to-sum}
\end{equation}
Taking imaginary and real parts proves the following pair.

\begin{theorem}[Thue--Morse product-to-sum identities]
For every integer \(m\ge0\) and complex \(z\),
\begin{equation}
  \boxed{
  \prod_{r=0}^{m}\sin\frac{z}{2^r}
   =\frac1{2^m}\sum_{n=0}^{2^m-1}\eps_n
     \sin\!\left(
       \frac{\pi m}{2}+\frac{(2n+1)z}{2^m}
     \right).}
  \label{eq:sine-product-to-sum}
\end{equation}
Also,
\begin{equation}
  \boxed{
  \cos z\prod_{r=1}^{m}\sin\frac{z}{2^r}
   =\frac1{2^m}\sum_{n=0}^{2^m-1}\eps_n
     \cos\!\left(
       \frac{\pi m}{2}+\frac{(2n+1)z}{2^m}
     \right).}
  \label{eq:cosine-product-to-sum}
\end{equation}
The empty product in \eqref{eq:cosine-product-to-sum} is \(1\) when \(m=0\).
\end{theorem}

Identity \eqref{eq:sine-product-to-sum} was posed in
\cite{reshetnikovTrigProducts}.  The complex derivation shows that it is
precisely a rescaled finite Fourier transform, and it supplies the cosine
companion for free.

\subsection{The discrete Fourier transform of a dyadic block}

For \(0\le k<2^m\), define
\begin{equation}
  \widehat\eps_m(k)=
  \sum_{n=0}^{2^m-1}\eps_n
    \e^{-2\pi\ii kn/2^m}.
  \label{eq:DFT-definition}
\end{equation}
Direct evaluation of the finite product gives
\begin{equation}
  \widehat\eps_m(k)=
  \prod_{j=0}^{m-1}
    \left(1-\e^{-2\pi\ii k2^j/2^m}\right).
  \label{eq:DFT-product}
\end{equation}
The last factor vanishes whenever \(k\) is even.  For odd \(k\), use
\(1-\e^{-\ii u}=2\ii\e^{-\ii u/2}\sin(u/2)\).

\begin{theorem}[Exact dyadic DFT]
For \(0\le k<2^m\),
\begin{equation}
  \boxed{
  \widehat\eps_m(k)=
  \begin{cases}
    0,&k\text{ even},\\[0.4em]
    (2\ii)^m\e^{-\pi\ii k(1-2^{-m})}
      \displaystyle\prod_{r=1}^{m}
      \sin\frac{\pi k}{2^r},&k\text{ odd}.
  \end{cases}}
  \label{eq:exact-DFT}
\end{equation}
\end{theorem}

Fourier inversion now gives a formula for each sign involving only odd
frequencies:
\begin{equation}
  \boxed{
  \eps_n=\frac1{2^m}
  \sum_{\substack{0\le k<2^m\\k\text{ odd}}}
    \widehat\eps_m(k)\e^{2\pi\ii kn/2^m},
  \qquad 0\le n<2^m,}
  \label{eq:DFT-inversion}
\end{equation}
with \(\widehat\eps_m(k)\) explicitly given by
\eqref{eq:exact-DFT}.  This is far less efficient than the binary-weight
formula, but it is the natural representation for spectral questions.

\section{Walsh formulas and the Riesz-product measure}
\label{sec:walsh-riesz}

\subsection{A one-point Walsh spectrum}
\label{sec:walsh}

For \(m\)-bit vectors \(a,n\in\mathbb F_2^m\), write
\(\langle a,n\rangle=\sum_ja_jn_j\pmod2\).  Since
\(\eps_n=(-1)^{\langle\bm1,n\rangle}\), where
\(\bm1=(1,\ldots,1)\), orthogonality of characters gives
\begin{equation}
  \boxed{
  \sum_{n<2^m}\eps_n(-1)^{\langle a,n\rangle}
   =2^m\ind_{a=\bm1}.}
  \label{eq:Walsh-spike}
\end{equation}
Thus a dyadic Thue--Morse block is spectrally maximally simple in the Walsh
basis: its transform has one nonzero coefficient.  The broad ordinary Fourier
spectrum in \eqref{eq:exact-DFT} is entirely a basis-change phenomenon.

The inverse character formula is correspondingly terse:
\begin{equation}
  \eps_n=(-1)^{\langle\bm1,n\rangle}.
  \label{eq:Walsh-character}
\end{equation}
Although this looks like the digit-sum definition, it changes the structural
viewpoint: \(\eps\) is a group character, not merely a recursively generated
word.

\subsection{Finite Riesz densities}
\label{sec:riesz}

The ordinary Fourier polynomial \(P_m\) produces the nonnegative density
\begin{equation}
  \rho_m(x)=\frac1{2^m}
    \left|P_m(\e^{2\pi\ii x})\right|^2.
  \label{eq:rho-definition}
\end{equation}
By \eqref{eq:finite-generating-product},
\begin{align}
  \rho_m(x)
  &=\prod_{j=0}^{m-1}
    \frac{|1-\e^{2\pi\ii2^jx}|^2}{2}
   =\boxed{
    \prod_{j=0}^{m-1}igl(1-\cos(2\pi2^jx)\bigr).}
  \label{eq:finite-Riesz-product}
\end{align}
The integral of \(\rho_m\) over \([0,1]\) is \(1\), by Parseval or by the
constant coefficient of \(|P_m|^2\).  Hence \(\rho_m(x)\dd x\) is a
probability measure.  These measures converge weakly to the classical
Thue--Morse measure, a purely singular continuous probability measure with
infinite Riesz-product representation
\begin{equation}
  \dd\mu_{\mathrm{TM}}(x)
   =\prod_{j\ge0}igl(1-\cos(2\pi2^jx)\bigr)\dd x,
  \label{eq:infinite-Riesz-product}
\end{equation}
where the displayed expression is interpreted as weak convergence of the
finite products, not as pointwise multiplication of ordinary densities
\cite{queffelec,baakeScaling}.

\subsection{Autocorrelation recursion}

Let
\begin{equation}
  \eta(r)=\lim_{N\to\infty}
    \frac1N\sum_{n=0}^{N-1}\eps_n\eps_{n+r},
  \label{eq:autocorrelation-definition}
\end{equation}
whose existence follows from unique ergodicity of the substitution system.
Splitting the average into even and odd indices and using
\eqref{eq:eps-recursion} gives
\begin{equation}
  \boxed{
  \eta(0)=1,
  \qquad
  \eta(2r)=\eta(r),
  \qquad
  \eta(2r+1)=-\frac{\eta(r)+\eta(r+1)}2.}
  \label{eq:autocorrelation-recursion}
\end{equation}
In particular,
\begin{equation}
  \eta(1)=-\frac{1+\eta(1)}2=-\frac13.
  \label{eq:eta-one}
\end{equation}
The numbers \(\eta(r)\) are the Fourier coefficients of
\(\mu_{\mathrm{TM}}\).  Thus the substitution recursion, the Riesz product,
and the spectral measure are three exact descriptions of the same object.

\section{Infinite series, integral transforms, and products}
\label{sec:infinite}

\subsection{Base-\(b\) Thue--Morse constants}

For an integer \(b\ge2\), define the base-\(b\) real number whose digits are
the Thue--Morse bits by
\begin{equation}
  \Theta_b=\sum_{n\ge0}\frac{\tau(n)}{b^{n+1}}.
  \label{eq:Theta-b-definition}
\end{equation}
Substituting \(z=1/b\) into \eqref{eq:T-from-E} yields the product formula
\begin{equation}
  \boxed{
  \Theta_b=\frac{1}{2(b-1)}
   -\frac{1}{2b}
    \prod_{j\ge0}\left(1-b^{-2^j}\right).}
  \label{eq:Theta-b-product}
\end{equation}
For \(b=2\), this is the classical Thue--Morse constant.  The signed companion
is even shorter:
\begin{equation}
  \sum_{n\ge0}\frac{\eps_n}{b^{n+1}}
  =\frac1b\prod_{j\ge0}\left(1-b^{-2^j}\right).
  \label{eq:signed-base-b}
\end{equation}

\subsection{A Dirichlet integral and its dyadic equation}

For \(a>0\), the bounded partial sums in
\eqref{eq:prefix-discrepancy} imply convergence of
\begin{equation}
  H(a)=\sum_{n\ge0}\frac{\eps_n}{n+a}.
  \label{eq:H-definition}
\end{equation}
Abel summation, or termwise integration first at radius \(r<1\) followed by
\(r\uparrow1\), gives
\begin{equation}
  H(a)=\int_0^1x^{a-1}E(x)\dd x
      =\int_0^1x^{a-1}
       \prod_{j\ge0}(1-x^{2^j})\dd x.
  \label{eq:H-integral}
\end{equation}
Splitting even and odd indices and using
\eqref{eq:eps-recursion} gives a two-scale functional equation:
\begin{equation}
  \boxed{
  H(a)=\frac12H\!\left(\frac a2\right)
       -\frac12H\!\left(\frac{a+1}{2}\right).}
  \label{eq:H-functional}
\end{equation}
This is the additive analogue of the multiplicative product equation below.

\subsection{A master two-parameter product}

For \(a,b>0\), define the product by dyadic blocks,
\begin{equation}
  G(a,b)=\lim_{m\to\infty}
   \prod_{n=0}^{2^m-1}
   \left(\frac{n+a}{n+b}\right)^{\eps_n}.
  \label{eq:G-definition}
\end{equation}
The logarithm converges by summation by parts because
\(\sum_{n<N}\eps_n\) is bounded and
\(\log((n+a)/(n+b))=\bigO(1/n)\) with a summable first difference.
The elementary laws
\begin{equation}
  G(a,a)=1,
  \qquad
  G(a,b)G(b,c)=G(a,c),
  \qquad
  G(b,a)=G(a,b)^{-1}
  \label{eq:G-cocycle}
\end{equation}
follow directly.

Pairing even and odd indices in each dyadic block gives the central functional
equation of the Thue--Morse gamma-product theory:
\begin{equation}
  \boxed{
  G(a,b)=
  \frac{G(a/2,b/2)}
       {G((a+1)/2,(b+1)/2)}.}
  \label{eq:G-functional}
\end{equation}
Indeed,
\begin{align*}
 &\left(\frac{2n+a}{2n+b}\right)^{\eps_{2n}}
 \left(\frac{2n+1+a}{2n+1+b}\right)^{\eps_{2n+1}}\\
 &\qquad=
 \left[
  \frac{(n+a/2)/(n+b/2)}
       {(n+(a+1)/2)/(n+(b+1)/2)}
 \right]^{\eps_n}.
\end{align*}
A broad family of closed evaluations of \(G\), including gamma-function
expressions and products with exponents \(\tau(n)\) rather than \(\eps_n\),
is developed in \cite{alloucheRiasatShallit}.

\subsection{The Woods--Robbins product}

The best-known evaluation is
\begin{equation}
  \boxed{
  \prod_{n\ge0}
   \left(\frac{2n+1}{2n+2}\right)^{\eps_n}
   =\frac1{\sqrt2}.}
  \label{eq:woods-robbins}
\end{equation}
It is \(G(1/2,1)\), but an elegant proof uses only index splitting
\cite{woods,robbins,alloucheShallitUbiquitous}.  Put
\begin{equation}
  A=\prod_{n\ge0}
    \left(\frac{2n+1}{2n+2}\right)^{\eps_n},
  \qquad
  B=\prod_{n\ge1}
    \left(\frac{2n}{2n+1}\right)^{\eps_n}.
  \label{eq:A-B-products}
\end{equation}
Then
\begin{align}
  AB
   &=\frac12\prod_{n\ge1}
      \left(\frac{n}{n+1}\right)^{\eps_n}
      \notag\\
   &=\frac12
      \prod_{j\ge0}
       \left(\frac{2j+1}{2j+2}\right)^{\eps_{2j+1}}
      \prod_{j\ge1}
       \left(\frac{2j}{2j+1}\right)^{\eps_{2j}}
      \notag\\
   &=\frac12A^{-1}B.
  \label{eq:woods-robbins-proof}
\end{align}
Cancel \(B\) to obtain \(A^2=1/2\); positivity gives
\eqref{eq:woods-robbins}.

\subsection{Three dyadic block-product limits}

The same even--odd splitting resolves three limits proposed in
\cite{reshetnikovProductLimits}.  With dyadic truncation,
\begin{align}
  \lim_{m\to\infty}
   \prod_{k=0}^{2^m-1}\left(k+\tfrac12\right)^{\eps_k}
   &=\frac12,
  \label{eq:block-product-one}\\
  \lim_{m\to\infty}
   \prod_{k=0}^{2^m-1}(k+1)^{\eps_k}
   &=\frac1{\sqrt2},
  \label{eq:block-product-two}\\
  \lim_{m\to\infty}
   \prod_{k=0}^{2^m-1}(k+1)^{(-1)^k\eps_k}
   &=\frac1{2\sqrt2}.
  \label{eq:block-product-three}
\end{align}
For example, pairing adjacent factors in the first two products reduces them
to \(G(1/4,3/4)\) and \(G(1/2,1)\), respectively:
\begin{align}
 \prod_{k<2^m}\left(k+\tfrac12\right)^{\eps_k}
  &=\prod_{j<2^{m-1}}
    \left(\frac{j+\tfrac14}{j+\tfrac34}\right)^{\eps_j},
  \label{eq:block-one-reduction}\\
 \prod_{k<2^m}(k+1)^{\eps_k}
  &=\prod_{j<2^{m-1}}
    \left(\frac{j+\tfrac12}{j+1}\right)^{\eps_j}.
  \label{eq:block-two-reduction}
\end{align}
The evaluations
\(G(1/4,3/4)=1/2\) and \(G(1/2,1)=1/\sqrt2\) are special cases of the
product theory in \cite{alloucheRiasatShallit}; the second is also
\eqref{eq:woods-robbins}.  A second split shows that the third limit is the
product of the first two.  These formulas illustrate why dyadic truncation is
the natural summation convention: it respects the exact block recursion.

\section{Iterated prefix sums and the bridge back to Fabius}
\label{sec:prefix-fabius}

Define inclusive discrete antiderivatives by
\begin{equation}
  S^{(0)}_n=\eps_n,
  \qquad
  S^{(k+1)}_n=\sum_{j=0}^{n}S^{(k)}_j.
  \label{eq:iterated-prefix-definition}
\end{equation}
The first prefix is already explicit from
\eqref{eq:prefix-discrepancy}:
\begin{equation}
  S^{(1)}_{2m}=\eps_m,
  \qquad
  S^{(1)}_{2m+1}=0.
  \label{eq:first-prefix-explicit}
\end{equation}
Repeated summation has a binomial-kernel formula.

\begin{theorem}[Binomial convolution for discrete antiderivatives]
For \(k\ge1\),
\begin{equation}
  \boxed{
  S^{(k)}_n=
  \sum_{j=0}^{n}
    \binom{n-j+k-1}{k-1}\eps_j.}
  \label{eq:prefix-binomial-convolution}
\end{equation}
Its generating function is
\begin{equation}
  \boxed{
  \sum_{n\ge0}S^{(k)}_nz^n
   =\frac{E(z)}{(1-z)^k}.}
  \label{eq:prefix-generating}
\end{equation}
\end{theorem}

\begin{proof}
One inclusive prefix multiplies an ordinary generating function by
\((1-z)^{-1}\).  After \(k\) prefixes one gets
\eqref{eq:prefix-generating}.  Expanding
\((1-z)^{-k}=\sum_{r\ge0}\binom{r+k-1}{k-1}z^r\) and taking a Cauchy product
gives \eqref{eq:prefix-binomial-convolution}.
\end{proof}

These formulas are formalized in \lean{ThueMorsePrefix.lean} and
\lean{ThueMorseGenerating.lean}.  The finite product also explains an exact
zero run at the end of every sufficiently large dyadic block.

\begin{theorem}[Dyadic terminal zero run]
If \(1\le k\le r\), then
\begin{equation}
  S^{(k)}_{2^r-k+i}=0,
  \qquad 0\le i<k,
  \label{eq:prefix-zero-run}
\end{equation}
while the entry immediately before the run and the next-block entry are
\begin{equation}
  S^{(k)}_{2^r-k-1}=(-1)^{r-k},
  \qquad
  S^{(k)}_{2^r}=-1.
  \label{eq:prefix-boundaries}
\end{equation}
\end{theorem}

\begin{proofidea}
Below degree \(2^r\), replace \(E(z)\) in
\eqref{eq:prefix-generating} by
\(P_r(z)=\prod_{j<r}(1-z^{2^j})\).  Since every factor contains \(1-z\),
\[
  \frac{P_r(z)}{(1-z)^k}
   =(1-z)^{r-k}
     \prod_{j=0}^{r-1}(1+z+\cdots+z^{2^j-1})

a polynomial of degree \(2^r-k-1\).  Therefore its next \(k\) coefficients,
at degrees \(2^r-k,\ldots,2^r-1\), vanish.  Comparing leading coefficients
gives the first value in \eqref{eq:prefix-boundaries}; the coefficient at
\(2^r\) follows from the first omitted factor
\(1-z^{2^r}\) or from the prefix recurrence.
\end{proofidea}

\subsection{Why this enters the Fabius development}

The Fabius connection uses exactly the chain now visible:
\begin{equation}
  \text{binary character}
  \Longrightarrow
  \prod_{j<r}(1-z^{2^j})
  \Longrightarrow
  \frac{\prod_{j<r}(1-z^{2^j})}{(1-z)^k}
  \Longrightarrow
  \text{iterated-prefix coefficients}.
  \label{eq:fabius-chain}
\end{equation}
After normalization by \(2^{\binom{k}{2}}\), suitable coefficient grids and
polygonal interpolants converge to the Fabius function.  The Prouhet zero of
order \(k\) is what creates the high-order flatness needed for the continuous
limit, while the finite Fourier product becomes the corresponding product of
sinc factors.  Thus the combinatorial and analytic appearances of
Thue--Morse in the Fabius theory are two specializations of the same finite
product, not separate accidents.

The present report deliberately stops before re-proving the quantitative
convergence results: those belong to the dedicated prefix-approximation
modules and companion convergence report.  What matters here is the formula
pipeline and the exact algebraic bridge.
